\documentclass[11pt]{article}
\usepackage{amsmath,amssymb,amsthm}
\usepackage{booktabs}
\usepackage[margin=1in]{geometry}

\newtheorem{theorem}{Theorem}
\newtheorem{lemma}[theorem]{Lemma}
\newtheorem{corollary}[theorem]{Corollary}

\theoremstyle{definition}
\newtheorem{definition}[theorem]{Definition}

\title{Problem 745: Sum of Squares II}
\author{Project Euler Solutions}
\date{}

\begin{document}
\maketitle

\section{Problem Statement}

For a positive integer~$n$, define $g(n)$ to be the maximum perfect square dividing~$n$. Define $S(N) = \sum_{n=1}^{N} g(n)$. Given $S(10)=24$ and $S(100)=767$, find $S(10^{14}) \bmod 1\,000\,000\,007$.

\section{Mathematical Foundation}

\begin{theorem}[Unique Squarefree Decomposition]
Every positive integer $n$ can be written uniquely as $n = a^2 b$ with $b$ squarefree, and $g(n) = a^2$.
\end{theorem}

\begin{proof}
Write $n = \prod p_i^{e_i}$. Set $a = \prod p_i^{\lfloor e_i/2\rfloor}$ and $b = \prod p_i^{e_i \bmod 2}$. Then $n = a^2 b$, $b$ is squarefree, and $a^2$ is the largest perfect square dividing~$n$. Uniqueness follows from unique prime factorization.
\end{proof}

\begin{theorem}[Jordan Totient Reformulation]
Define the Jordan totient function of order~$2$:
\[
J_2(e) = \sum_{a \mid e} a^2\,\mu(e/a) = e^2 \prod_{p \mid e}\!\left(1 - \frac{1}{p^2}\right).
\]
Then:
\[
S(N) = \sum_{e=1}^{\lfloor\sqrt{N}\rfloor} J_2(e)\left\lfloor\frac{N}{e^2}\right\rfloor.
\]
\end{theorem}

\begin{proof}
From the squarefree decomposition:
\[
S(N) = \sum_{a=1}^{\lfloor\sqrt{N}\rfloor} a^2 \sum_{\substack{b \le N/a^2 \\ b\text{ sqfree}}} 1 = \sum_{a=1}^{\lfloor\sqrt{N}\rfloor} a^2 \sum_{d=1}^{\lfloor\sqrt{N/a^2}\rfloor} \mu(d)\left\lfloor\frac{N}{a^2 d^2}\right\rfloor.
\]
Substituting $e = ad$:
\[
S(N) = \sum_{e=1}^{\lfloor\sqrt{N}\rfloor}\left\lfloor\frac{N}{e^2}\right\rfloor \underbrace{\sum_{a \mid e} a^2\,\mu(e/a)}_{= J_2(e)}.
\]
The multiplicative formula $J_2(e) = e^2\prod_{p \mid e}(1-1/p^2)$ follows from evaluating $J_2$ at prime powers: $J_2(p^k) = p^{2k} - p^{2(k-1)}$.
\end{proof}

\begin{lemma}[Sieve for $J_2$]
The values $J_2(e)$ for $1 \le e \le L$ are computed by initializing $J_2(e) = e^2$ and, for each prime~$p \le L$, multiplying $J_2(m)$ by $(p^2-1)/p^2$ for every multiple~$m$ of~$p$. This runs in $O(L\log\log L)$ time.
\end{lemma}

\begin{proof}
The sieve correctly accumulates $\prod_{p \mid e}(1-1/p^2)$ for each~$e$. The time bound follows from the standard prime-harmonic analysis.
\end{proof}

\section{Algorithm}

\begin{enumerate}
\item Set $L = \lfloor\sqrt{N}\rfloor = 10^7$.
\item Sieve $J_2(e) \bmod p$ for $1 \le e \le L$:
  \begin{itemize}
  \item Initialize $J_2(e) = e^2 \bmod p$.
  \item For each prime $q \le L$, for each multiple~$m$: $J_2(m) \leftarrow J_2(m) \cdot (q^2-1) \cdot (q^2)^{-1} \bmod p$.
  \end{itemize}
\item Compute $S = \sum_{e=1}^{L} J_2(e) \cdot \lfloor N/e^2\rfloor \bmod p$.
\end{enumerate}

\section{Complexity Analysis}

\begin{itemize}
\item \textbf{Time:} $O(\sqrt{N}\log\log\sqrt{N})$ for the sieve, $O(\sqrt{N})$ for summation.
\item \textbf{Space:} $O(\sqrt{N})$ for the $J_2$ array and prime sieve.
\end{itemize}

\section{Answer}

\[
\boxed{94586478}
\]

\end{document}
